%lezione 7 del 23 marzo 2026
\section{Postulato 2 della Meccanica Quantistica}
\noindent\fbox{%
    \parbox{\textwidth}{%
    \textbf{La misura con precisione assoluta di tutte le osservabili che compongono un SCOC mette il
    sistema in un loro autostato simultaneo di dimensione 1.}
    }%
}
\subsection{Il collasso della funzione d'onda}

Questo postulato è la versione più generale della definizione di collasso della funzione d'onda:
nel caso di SCOC composto da una sola osservabile $A$, significa che la misura di $A$ \textbf{modifica lo
stato del sistema e lo pone nell'autostato normalizzato di} $\mathbf{A}$ \textbf{associato all'autovalore
misurato}
$$\ket{\psi}\rightarrow \frac{P_{a_i}\ket{\psi}}{\left|\left| P_{a_i}\ket{\psi} \right|\right|}$$
cioè se, dopo la prima misura, si misura nuovamente l'osservabile $A$ si ottiene \textbf{SEMPRE} lo stesso
risultato $a_i$.


\textbf{Esempio} Sia dato un sistema fisico con SCOC $\left\{A,B\right\}$ e lo stato 
$$\ket{\psi} = \sum_{k,n} c_{kn}\ket{a_k,b_n}$$
Se una misura di $A$ restituisce il valore $a_i$, lo stato diventa (a meno di normalizzazioni)
$$\ket{\psi}\rightarrow P_{a_i}\ket{\psi}=
\sum_{k,n}c_{kn}\sum_{m}\ket{a_i,b_m}\bra{a_i,b_m}\ket{a_k,b_n} = \sum_{m}c_{im}\ket{a_i,b_m}:=\ket{\psi'}
$$
dove si è utilizzato il fatto che $\bra{a_i,b_m}\ket{a_k,b_n} = \delta_{ik}\delta_{mn}$.

Se poi su questo stato si misura $B$ e la misura restituisce il valore $b_l$, lo stato diventa
$$\ket{\psi'} \rightarrow P_{b_l}\ket{\psi'} =
\sum_{m}c_{im}\sum_{j}\ket{a_j,b_l}\bra{a_j,b_l}\ket{a_i,b_m} = c_{il}\ket{a_i,b_l}$$
che è, a meno di normalizzazioni, un autostato simultaneo del SCOC, elemento dell'autospazio comune di 
dimensione 1.

Tutto questo ragionamento è analogo per gli stati misti, con la differenza che la misura di $A$ su $\rho$
comporta la trasformazione
$$\rho\rightarrow \frac{P_{a_i}\rho P_{a_i}}{Tr\left(P_{a_i}\rho P_{a_i}\right)}:=\rho'$$
che mantiene la "normalizzazione" di $\rho$, cioè $Tr(\rho')=1$

\textbf{Esempio, stato misto} Si consideri la stessa situazione dell'esempio precedente, dato però lo
stato misto
$$\rho = \sum_{\alpha}\omega_\alpha\ket{\psi_\alpha}\bra{\psi_\alpha}$$
Dopo una misura di $A$ che restituisce $a_i$ lo stato diventa
\begin{align*}
    \rho\rightarrow P_{a_i}\rho P_{a_i}&=
    \sum_{\alpha}\omega_\alpha\sum_{m,n}\ket{a_i,b_m}\bra{a_i,b_m}\ket{\psi_\alpha}\bra{\psi_\alpha}
    \ket{a_i,b_n}\bra{a_i,b_n}=\\
    &=\sum_{\alpha}\omega_\alpha\sum_{m,n} c_{im}^{(\alpha)}c_{in}^{(\alpha)*}\ket{a_i,b_m}\bra{a_i,b_m}:=
    \rho'
\end{align*}
E se, similmente a prima, una successiva misura di $B$ restituisce $b_l$ si trova
\begin{align*}
    \rho'\rightarrow P_{b_l}\rho'P_{b_l} &=
    \sum_{\alpha}\omega_\alpha\sum_{m,n} c_{im}^{(\alpha)}c_{in}^{(\alpha)*}\sum_{k,j}\ket{a_j,b_l}
    \bra{a_j,b_l}\ket{a_i,b_m}\bra{a_i,b_n}\ket{a_k,b_l}\bra{a_k,b_l} =\\
    &= \sum_{\alpha} |c_{il}^{(\alpha)}|^2\ket{a_i,b_l}\bra{a_i,b_l}
    := d_{il}\ket{a_i,b_l}\bra{a_i,b_l}
\end{align*}
che è, a meno di normalizzazioni, uno stato puro che è autostato simultaneo del SCOC.

\textbf{Oss.} Il postulato 2 fornisce sia un modo di costruire autostati di operatori, come già visto, sia
un metodo per costruire stati misti.\\
Dato un sistema con SCOC 1-dimensionale $\{A\}$,
si preparino $N$ stati fisici identici 
$$\ket{\psi} = \sum_{i}^{}c_i\ket{a_i}$$
e si effettui su uno di essi una misura di $A$:
$$\ket{\psi}\rightarrow \frac{c_i\ket{a_i}}{||c_i\ket{a_i}||} = e^{i\varphi_i}\ket{a_i}\sim\ket{a_i}$$
la probabilità di misurare $a_i$ è $\mathcal{P}_{\psi}(a_i)=|c_i|^2:=\omega_i$: se si svolge questo
procedimento su tutti i sistemi e si considera la sovrapposizione dei risultati, si trova lo stato misto
$$\rho = \sum_{k}|c_k|^2\ket{a_k}\bra{a_k}$$
Questo procedimento costruisce uno stato incoerente proprio perché le \textit{coerenze} dello stato
iniziale sono date dalle fasi dei $c_i$, che sono state cancellate da questo procedimento.

\textbf{Problema} Come si può "misurare con precisione assoluta"?\\
In meccanica classica, questo non è \textit{mai} possibile, in quanto i valori di ogni osservabile sono
continui.\\
In meccanica quantistica è invece possibile \textit{solo se} l'osservabile misurata ha
\textbf{spettro discreto} e l'apparato sperimentale ha risoluzione minore della distanza tra ciascuna 
coppia di autovalori consecutivi.\\
Per le osservabili con spettro continuo non è invece possibile, perché la risoluzione sperimentale non è
mai abbastanza fine da "costringere" il sistema esattamente in un solo autostato.

\textbf{Oss.} Questo è il motivo profondo per cui gli autostati dello spettro continuo di un'osservabile
non sono elementi dello spazio di Hilbert del sistema, ma sono comunque basi generalizzate di esso:
se una misura di posizione in 1D restituisce $x=(4.0\pm0.1)$fm, lo stato del sistema è una sovrapposizione
continua di autostati di $X$ ed è normalizzabile; viceversa, se la precisione fosse infinita, lo stato non
sarebbe normalizzabile e non si potrebbe definire una distribuzione di probabilità.
%\newpage%?
\section{Postulato 3 della Meccanica Quantistica}
Si enuncia infine il postulato dell'evoluzione temporale.

\noindent\fbox{%
    \parbox{\textwidth}{%
	\textbf{Dato un sistema fisico con Hamiltoniana H e un tempo iniziale $\mathbf{t_0}$,
	l'operatore evoluzione temporale è l'unico operatore $\mathbf{U(t_0,t)}$ soluzione
	dell'equazione differenziale di Scrödinger}
	$$i\hbar\dv[]{}{t}U(t,t_0) = H(t)U(t,t_0)$$
	\textbf{con condizione iniziale} $U(t_0,t_0) = \mathbb{I}$.\\
	\textbf{Dato uno stato $\mathbf{\ket{\psi}}$ a $\mathbf{t=t_0}$, il valore medio di
	un'osservabile F al tempo $\mathbf{t>t_0}$ è}
	$$\langle F\rangle_\psi(t):=\bra{\psi}U^{\dagger}(t,t_0)FU(t,t_0)\ket{\psi}$$
    }%
}

L'operatore $U(t,t_0)$ è quindi l'ultimo tassello del puzzle: come esplicitato nel postulato, permette
di fare predizioni sul valore medio delle osservabili nel tempo.

\textbf{Oss.} $U(t,t_0)$ è unitario: si considerino infatti l'equazione di Schrödinger e la sua hermitiana
coniugata
\begin{align*}
    i\hbar\dv[]{U}{t} = HU &\rightarrow U^{\dagger}i\hbar\dv[]{U}{t} = U^{\dagger}HU\\
    -i\hbar\dv[]{U^{\dagger}}{t} = U^{\dagger}H &\rightarrow -i\hbar\dv[]{U^{\dagger}}{t}U = U^{\dagger}HU
\end{align*}
Sottraendo membro a membro si ottiene
$$ i\hbar\left(U^{\dagger}\dv[]{U}{t}+\dv[]{U^{\dagger}}{t}\right) = 0 \leftrightarrow
\dv[]{}{t}\left(U^{\dagger}U\right) = 0 \leftrightarrow U^{\dagger}U=\text{cost.}$$
Ricordando che $U(t_0,t_0) = \mathbb{I} = U^{\dagger}(t_0,t_0)$, si ha che
$$ U(t,t_0)U^{\dagger}(t,t_0)=\mathbb{I}\mkern9mu \forall t$$
Se $\pdv[]{H}{t}=0$, la forma esplicita di $U(t,t_0)$ è facile da ricavare:
$$\dv[]{U}{t} = -\frac{i}{\hbar}HU \leftrightarrow \int_{U(t_0,t_0)}^{U(t_0,t)}\frac{dU}{U} =
-\frac{i}{\hbar}H\int_{t_0}^{t}d\tau \leftrightarrow U(t,t_0)=\exp\left\{-\frac{i}{\hbar}H(t-t_0)\right\}$$
Dove, dato un operatore $A$,
$$e^A:=\sum_{n=0}^{\infty}\frac{1}{n!}A^n$$

Ci sono due schemi interpretativi dell'evoluzione temporale del valor medio di un'osservabile che
però \textbf{non} portano a conseguenze fisiche diverse, in quanto l'unica parte della teoria che ha
significato fisico è la previsione dei valori medi.
\subsection{Schema di Scrödinger}
Nello schema di Schrödinger si utilizza il concetto di \textit{stato dipendente dal tempo}
\begin{gather*}
\ket{\psi(t)}:=U(t,t_0)\ket{\psi(t_0)}\\
\rho(t)=\sum_{k}^{}\omega_k\ket{\psi_k(t)}\bra{\psi_k(t)} = U(t,t_0)\rho(t_0)U^{\dagger}(t,t_0)
\end{gather*}
Poiché $UU^{\dagger}=\mathbb{I}$, $\bra{\psi(t)}\ket{\psi(t)} = \bra{\psi(t_0)}\ket{\psi(t_0)}$.

Gli autostati degli operatori \textit{non} evolvono nel tempo, ma i coefficienti di Fourier sì:
$$A\ket{a_i}=a_i\ket{a_i},\quad\quad\ket{\psi(t)} = \sum_{i}\ket{a_i}\bra{a_i}\ket{\psi(t)} =
\sum_{i}c_i(t)\ket{a_i}$$

\subsection{Schema di Heisenberg}
Nello schema di Heiseberg, invece, l'evoluzione temporale è associata alle osservabili
$$A(t):=U^\dagger(t,t_0) A(t_0)U(t,t_0)$$

Gli autostati al contrario di prima evolvono nel tempo
$$A(t_0)\ket{a(t_0)}=a\ket{a(t_0)}\rightarrow
U^\dagger A(t_0)UU^\dagger\ket{a(t_0)}=a U^\dagger\ket{a(t_0)}\rightarrow
A(t)\ket{a(t)} = a\ket{a(t)}$$
dove si è definito $\ket{a(t)}:=U^\dagger(t,t_0)\ket{a(t_0)}$.
\subsection{Costanti del moto}
Similmente alla meccanica Hamiltoniana, dove se $\pdv[]{H}{t}=0$,
$$\left\{H,F\right\}_{PB}=0\leftrightarrow\dv[]{F}{t}=0$$
anche in meccanica quantistica vale che, se $H$ è indipendente dal tempo,
$$\bra{\psi}e^{\frac{i}{\hbar}H(t-t_0)} Fe^{-\frac{i}{\hbar}H(t-t_0)}\ket{\psi} = 
\bra{\psi}FU^\dagger U\ket{\psi} = \bra{\psi}F\ket{\psi} \leftrightarrow
\dv[]{}{t}\langle F\rangle_\psi(t)=0
$$
dove si è utilizzato il fatto che 
\begin{gather*}
[U,F]=\sum_{n=0}^{\infty}\left(-\frac{i}{\hbar}(t-t_0)\right)^n\frac{1}{n!}[H^n,F]
=\sum_{n=0}^{\infty}\left[\left(-\frac{i}{\hbar}(t-t_0)\right)^n\frac{1}{n!}
\left(\sum_{k=1}^{n-1}a_{k,n}H^{n-k-1}[H,F]H^{k}\right)\right]=0
\end{gather*}
se e solo se $[H,F] = 0$.

%{\large DA DIMOSTRARE RIVEDI RIVEDI}
%$$[A^n,B] = \sum_{k=0}^{n-1}\binom{n-1}{k}A^{n-1-k}[A,B]A^k$$
\textbf{Oss.} Gli autostati dell'energia conservano tutte le osservabili:
se $H$ è indipendente dal tempo, per gli autostati dell'energia si ha
$$H\ket{\psi}=E\ket{\psi} \rightarrow \ket{\psi(t)}=U(t,t_0)\ket{\psi} = e^{-\frac{i}{\hbar}E(t-t_0)}
\ket{\psi} \rightarrow \bra{\psi}U^\dagger FU\ket{\psi} = \bra{\psi}F\ket{\psi}\mkern9mu\forall t,F$$
